\chapter{Related Work}\label{chap:related}

\section{Generalized Additive Models and Interpretability}

\Glspl{GAM}, introduced by \citet{hastie1986generalized}, model the response as a sum of univariate functions.
\citet{lou2012intelligible, lou2013accurate} proposed \glspl{GA2M}, augmenting GAMs with pairwise interaction terms, and Microsoft's \gls{EBM} \citep{nori2019interpretml} brought these to production use.
More recently, \citet{agarwal2021neural} introduced Neural Additive Models (NAMs), which inherit the additive structure (and thus the identifiability problem).
\citet{chang2021how} showed that different GAM training algorithms produce substantially different shape functions for the same data---directly analogous to our finding that different weight specifications yield different purified effects.

\section{Functional ANOVA and Sensitivity Analysis}

The \gls{fANOVA} decomposition has roots in global sensitivity analysis.
\citet{sobol2001global} introduced Sobol' indices based on the orthogonal decomposition of a function under a product measure (assuming independent inputs).
\citet{hooker2007generalized} extended the framework to dependent variables, showing that the standard decomposition can misattribute interaction variance to main effects and vice versa when features are correlated---the phenomenon our experiments quantify.

\section{Feature Interactions and Their Detection}

\citet{friedman2008predictive} proposed the H-statistic for measuring interaction strength, conceptually related to the interaction variance share we compute after purification.
\citet{lundberg2017unified} introduced \gls{SHAP} interaction values, which decompose pairwise effects at the individual prediction level rather than as global shape functions, making them complementary to purification.
\citet{hooker2021unrestricted} and \citet{strobl2008conditional} showed that permutation-based importance measures are unreliable when features are dependent---paralleling our finding that the weighting distribution $w$ plays an analogous role to the reference distribution.

\section{Partial Dependence and Accumulated Local Effects}

\citet{friedman2001greedy} introduced partial dependence plots (PDPs), which implicitly assume feature independence and can mislead when features are correlated.
\citet{apley2020visualizing} developed accumulated local effects (ALE) plots as an alternative that avoids extrapolation by restricting attention to observed feature combinations.
Both face the same marginal-vs-conditional tension that our work addresses through weight specification.

\section{Purification and Identifiability}

\citet{lengerich2020purifying} identified the non-identifiability of additive-with-interaction models and proposed the mass-moving purification algorithm.
They proved convergence guarantees and demonstrated the algorithm on COMPAS, pneumonia risk, and recidivism datasets, showing that purification qualitatively changes which features appear important and how interactions are attributed.

However, their analysis focused on comparing purified versus unpurified effects, rather than studying sensitivity to the weighting distribution $w$.
This gap---understanding how much purified effects change when $w$ is varied---is the primary contribution of our work.

\section{Algorithmic Fairness and the German Credit Dataset}

The German Credit dataset \citep{dua2019uci} is a standard benchmark in fair lending research, containing sensitive attributes such as age and sex.
The COMPAS recidivism system \citep{angwin2016machine} highlighted how model interpretations can have profound societal consequences, with debates centering on how ``bias'' is measured and attributed.
The purification algorithm is directly relevant: by choosing different weighting distributions, a practitioner could shift variance between main effects involving sensitive attributes and their interactions, raising questions about the trustworthiness of purification-based audits (\cref{sec:ethical}).
